\begin{textsample}{POS Dim 1 – 2010\_c5 – Score 4.00 – 2010\_c5/enade\_2010\_farm\_c5f3fda3be.txt}  \label{ex:f1_pos_136}
CAPÍTULO 5 DISTRIBUIÇÃO DOS CONCEITOS 5.1 \textbf{PANORAMA} NACIONAL DA DISTRIBUIÇÃO DOS CONCEITOS A Tabela 5.1 apresenta a quantidade e distribuição de cursos de Farmácia participantes do ENADE/2010 por faixa de conceito e Grande Região.
Observando -se os dados da Tabela 5.1, nota -se que, dos 340 cursos participantes, 38,7 \% ( 132 cursos ) classificaram -se com conceito 3, o valor modal.
Este foi o conceito modal em todas as regiões brasileiras, com percentuais que variaram de 27,9 \%, na região Centro-Oeste a 43,8 \%, na Sudeste.
O conceito 4 foi o segundo mais frequente nas regiões Sudeste ( 15,4 \%, correspondendo a 26 cursos ), Sul ( 23,2 \%, relativos a 16 cursos ) e a Centro-Oeste ( 19,4 \%, com sete cursos, onde houve o mesmo percentual e número de cursos para o conceito 2 ).
O conceito 1 foi verificado nas regiões Sudeste ( 1,8 \%, relacionado a três cursos ) e Sul ( dois cursos ou 2,9 \% do total regional ).
O conceito 2, nas regiões Sudeste e Sul, foi a classificação, respectivamente, de 24 e de seis cursos.
Considerando -se o total nacional, houve, ainda, 23 cursos ( 6,8 \% ) que receberam conceito 5 e cinco cursos ( 1,5 \% ) que receberam conceito 1.
Dos 340 cursos, 76 ( 22,4 \% ) ficaram sem conceito ( SC ).
A região Norte participou com 19 cursos, 5,6 \% do total nacional.
Nesta região, o conceito 3 foi atribuído a seis dos 19 cursos participantes, equivalentes a 31,6 \% do total regional.
Um pouco menos de 1/3 dos cursos ( 31,6 \%, correspondendo a 6 cursos ) não receberam conceito 2 ( 21,1 \%, correspondendo a quatro cursos ) e conceito 4 ( 15,8 \%, correspondendo a três cursos ).
A região Nordeste participou com 47 cursos, 13,8 \% do total nacional.
Destes, 15 cursos, 31,9 \% em termos regionais, obtiveram conceito 3.
Os conceitos 2 e 4 foram atribuídos, respectivamente, a seis ( 12,8 \% ) e cinco ( 10,6 \% ) cursos.
O conceito 5 foi atribuído a três cursos ( 6,4 \% ).
Nessa região, 38,3 \% dos cursos ( 18 cursos ) ficaram sem conceito.
Dos 169 cursos participantes da região Sudeste, 74 ( 43,8 \% ) obtiveram conceito 3, o conceito modal.
O conceito 2 foi atribuído a 24 cursos ( 14,2 \% ) e o conceito 4, a 26 ( 15,4 \% ).
Três cursos ( 1,8 \% ) receberam conceito 1 nessa região.
Outros dez cursos ( 5,9 \% ) receberam conceito 5 e os 32 cursos \textbf{restantes} ficaram sem conceito.
Na região Sul, 27 cursos ( 39,2 \% ) dos 69 cursos participantes receberam o conceito modal, 3.
O conceito 1 foi atribuído a 2 cursos ( 2,9 \% ) e o conceito 2, a seis deles ( 8,7 \% ).
O conceito 4 foi a avaliação recebida por 16 cursos ( 23,2 \% ), e o conceito 5, a sete cursos ( 10,1 \% ).
Do total de cursos da região Sul, 11 ( 15,9 \% ) ficaram sem conceito.
Dos 36 cursos participantes na região Centro-Oeste, nove ( 25,0 \% ) ficaram sem conceito, sete receberam conceito 2 e dez receberam conceito 3 ( 19,4 \% e 27,9 \% ), o conceito modal.
Nos cursos \textbf{restantes}, o conceito 4 foi atribuído a sete deles ( 19,4 \% ), e o conceito 5, a três cursos ( 8,3 \% ). 5.2 CONCEITOS POR CATEGORIA ADMINISTRATIVA E POR GRANDE REGIÃO A Tabela 5.2 apresenta a distribuição dos cursos participantes do ENADE/2010 de Farmácia, por Categoria Administrativa, de acordo com os conceitos por eles alcançados, segundo as Grandes Regiões brasileiras.
Dos 340 cursos participantes, 61 ( 17,9 \% ) eram \textbf{ministrados} em instituições públicas e 279 ( 82,1 \% ), em privadas.
De acordo com as informações da Tabela 5.2, as instituições privadas foram as que apresentaram, em termos nacionais, a maior proporção de cursos com conceito 3, equivalente a 43,0 \% ( 120 cursos ) do total dessa categoria.
Dos 279 cursos participantes nesta Categoria Administrativa, o conceito 4 foi atribuído a 43 cursos.
Classificaram -se, ainda, 43 cursos com conceito 2 e quatro com conceito 5.
O conceito 1 foi atribuído a cinco cursos nesta categoria e 64 deles ficaram sem conceito.
Na rede pública, o conceito modal foi 5 ( com 19 cursos ).
Entre os demais cursos participantes, quatro receberam conceito 2, doze conceito 3 e 14 conceito 4.
Nesta categoria, 12 cursos ficaram sem conceito.
Na análise por região, observa -se que, na região Norte, as instituições privadas participaram com 14 cursos, dos quais 4 ficaram sem conceito.
O conceito modal para as instituições privadas na região Norte foi 3 ( com cinco cursos ).
Os demais foram avaliados com conceito 2 ( três cursos ) e conceito 4 ( dois cursos ).
As instituições públicas participaram com cinco cursos, dos quais dois ficaram sem conceito.
Os conceitos 2, 3 e 4 ficaram com um curso cada.
Na região Nordeste, a rede privada concentrou 32 dos 47 cursos participantes, correspondentes a 68,1 \% do total da região.
Nesta Categoria Administrativa na região Nordeste, predominaram os cursos com conceito 5 ( 15 cursos ), sendo este, portanto, o conceito modal.
Os demais foram avaliados com os seguintes conceitos : conceito 1 ( cinco cursos ), conceito 2 ( 11 cursos ) e conceito 3 ( um curso ).
Nenhum curso ficou sem conceito.
As instituições públicas, nesta região, participaram com 15 cursos ( 31,9 \% ), sendo que não houve, também, casos de cursos sem conceito.
O conceito modal para esta Categoria Administrativa ficou entre 2 e 3 ( quatro cursos cada ).
Os demais foram avaliados com os conceitos 1 ( um curso ), 4 ( três cursos ) e 5 ( também três cursos ).
Na região Sudeste, a proporção de cursos da rede privada, 89,3 \%, foi mais elevada do que nas demais regiões brasileiras, correspondendo a 151 dos 169 cursos participantes.
Nesta categoria na região, o conceito modal foi 3 ( 71 cursos ).
Os demais foram avaliados com conceito 1 ( três cursos ), conceito 2 ( 24 cursos ), conceito 4 ( 23 cursos ) e conceito 5 ( três cursos ).
A partir da combinação de Categoria Administrativa e Grande Região, observa -se que 27 cursos não receberam conceito.
Entre os 18 cursos em instituições públicas na região, a categoria modal foi a que recebeu conceito 5 ( sete cursos ).
Os demais foram avaliados com conceito 3 ( três cursos ) e conceito 4 ( três cursos ).
Nesta categoria na região Sudeste, cinco cursos ficaram sem conceito.
As instituições privadas concentraram 53 dos 69 cursos participantes da região Sul, correspondentes a 76,8 \% do total regional.
Destes, 23 ficaram classificados no conceito modal, 3.
Os demais foram avaliados com conceito 1 ( dois cursos ), conceito 2 ( cinco cursos ) e conceito 4 ( 13 cursos ).
Na combinação de Categoria Administrativa e Grande Região, dez cursos não receberam conceito.
As instituições públicas nesta região participaram com 16 cursos, dos quais sete receberam o conceito modal, correspondente a 5.
Os demais foram avaliados com conceito 2 ( um curso ), conceito 3 ( quatro cursos ), conceito 4 ( três cursos ), além de um curso que ficou sem conceito.
Na região Centro-Oeste, 29 dos 36 cursos participantes eram de instituições privadas ( 80,6 \% em termos regionais ).
Destes, dez concentraram -se no conceito modal, 3, e oito ficaram sem conceito.
Os demais receberam conceito 2 ( seis cursos ), conceito 4 ( quatro cursos ) e conceito 5 ( um curso ).
Na rede pública, dos sete cursos na região Centro-Oeste, três foram classificados no conceito modal, 4, e um ficou sem conceito.
Os demais foram avaliados com conceito 2 ( um curso ) e conceito 5 ( dois cursos ). 5.3 CONCEITOS POR ORGANIZAÇÃO ACADÊMICA E POR GRANDE REGIÃO Na Tabela 5.3 encontra -se a distribuição dos conceitos atribuídos aos cursos participantes do ENADE/2010 na Área de Farmácia, por Organização Acadêmica, segundo as Grandes Regiões brasileiras.
Dos 340 cursos participantes, 165 eram oferecidos em Universidades, 59 em Centros Universitários e os demais 116 em Faculdades.
De acordo com os dados apresentados, a maioria dos cursos avaliados com conceito 5 eram vinculados a Universidades ( 22 dos 23 cursos com este conceito ).
Esse tipo de Organização Acadêmica teve o conceito 3 como modal, com 60 cursos.
Além dos 25 cursos que ficaram sem conceito, os demais foram avaliados com conceito 1 ( quatro cursos ), conceito 2 ( 13 cursos ) e conceito 4 ( 41 cursos ).
Entre os cursos em Centros Universitários, o conceito modal também foi 3 ( com 26 cursos ), e dez ficaram sem conceito.
Os outros cursos neste tipo de Organização Acadêmica receberam os conceitos 1 ( um curso ), 2 ( 11 cursos ) e 4 ( também 11 cursos ).
Nas Faculdades, dos 116 cursos, 41 ficaram sem conceito e 46 receberam o mesmo conceito modal dos demais tipos de Organização Acadêmica, 3.
Os demais cursos receberam conceito 2 ( 23 cursos ), conceito 4 ( cinco cursos ) e conceito 5 ( apenas um curso ).
Considerando -se separadamente as regiões brasileiras, verifica -se que, na região Norte, as Universidades concentraram 5 dos 19 cursos participantes.
Dos cursos em Universidades, dois ficaram sem conceito e os três \textbf{restantes} foram avaliados com conceito 3 ( um curso ) e conceito 4 ( dois cursos – o conceito modal na combinação de Organização Acadêmica e Grande Região ).
Os Centros Universitários foram representadas por seis cursos, sendo um sem conceito.
Os demais cursos receberam conceito 2 ( três cursos – o conceito modal ), conceito 3 ( um curso ) e conceito 4 ( também um curso ).
As Faculdades participaram com oito cursos na região.
Destes, três ficaram sem conceito e os cinco demais receberam os conceitos 2 ( um curso ) e 3 ( 4 cursos – o conceito modal ).
Na região Nordeste, as Universidades participaram com 18 dos 47 cursos na Área de Farmácia da região.
Entre os cursos em Universidades, quatro ficaram sem conceito.
O conceito modal ficou distribuído entre os conceitos 3 e 4, ambos com cinco cursos.
Os demais cursos receberam conceito 2 ( um curso ) e conceito 5 ( três cursos ).
Os Centros Universitários \textbf{contaram} com três cursos participantes na região Nordeste.
Todos os cursos receberam conceito e o conceito modal foi 3 ( com dois cursos ).
O curso \textbf{restante} recebeu conceito 2.
As Faculdades foram representadas por 26 cursos, dos quais 14 não receberam conceito.
O conceito modal foi 3 ( com oito cursos ).
Os quatro cursos \textbf{restantes} receberam conceito 2.
Na região Sudeste as Universidades concentraram 82 dos 169 cursos da região.
Entre os cursos em Universidades, o conceito modal foi 3 ( com 35 cursos ) e 12 ficaram sem conceito.
Os demais cursos receberam conceito 1 ( dois cursos ), conceito 2 ( nove cursos ), conceito 4 ( 15 cursos ) e conceito 5 ( nove cursos ).
Os Centros Universitários participaram com 38 cursos na região Sudeste, dos quais 18 obtiveram o conceito modal, 3, e sete ficaram sem conceito.
Os demais receberam conceito 1 ( um curso ), conceito 2 ( quatro cursos ), e conceito 4 ( oito cursos ).
As Faculdades foram representadas por 49 cursos na região Sudeste, que se distribuíram nos conceitos 2 ( 11 cursos ), 3 ( 21 cursos, o conceito modal ), 4 ( três cursos ) e 5 ( um curso ), ou ficaram sem conceito ( 13 cursos ).
Dos 69 cursos da região Sul, 47 eram de Universidades, para os quais o conceito modal foi 3 ( com 17 cursos ).
Nesse tipo de organização, cinco cursos ficaram sem conceito e os demais receberam conceito 1 ( dois cursos ), conceito 2 ( dois cursos ), conceito 4 ( 14 cursos ) e conceito 5 ( sete cursos ).
Os Centros Universitários tiveram sete cursos participantes, sendo que um deles não recebeu conceito.
Os demais seis cursos foram distribuídos, em mesma quantidade ( dois cursos cada ), aos conceitos 2, 3 e 4.
Foram 15 os cursos vinculados a Faculdades, e destes, cinco ficaram sem conceito.
Os dez \textbf{restantes} receberam conceito 2 ( dois cursos ) e conceito 3 ( oito cursos, o conceito modal ).
Na região Centro-Oeste, 13 dos 36 cursos eram de Universidades.
Nesse tipo de organização, dois cursos ficaram sem conceito e o conceito modal foi 4 ( com cinco cursos ).
Os seis cursos \textbf{remanescentes} foram avaliados com conceito 2 ( um curso ), conceito 3 ( dois cursos ) e conceito 5 ( três cursos ).
Os Centros Universitários \textbf{contaram} com 5 cursos, um dos quais não recebeu conceito e três receberam conceito 3, o conceito modal.
O curso \textbf{restante} recebeu conceito 2.
Dos 18 cursos em Faculdades da região Centro-Oeste, seis ficaram sem conceito.
O conceito modal ficou entre 2 e 3, ambos com cinco cursos.
Os dois outros cursos foram avaliados com conceito 4.

% matched lemmas: contar, ministrar, panorama, remanescente, restante
\end{textsample}
